\documentclass[11pt]{article}

\usepackage[final]{acl}
\usepackage{times}
\renewcommand{\ttdefault}{lmtt}
\usepackage{latexsym}
\usepackage[T1]{fontenc}
\usepackage[utf8]{inputenc}
\usepackage{microtype}
\usepackage{amsmath}
\usepackage{amssymb}

\title{Contradiction-Aware Editing in Multi-Turn Settings\\for Graph Language Models}

\author{
Arihant Rastogi \\ 2023111012 \And
Prasoon Dev \\ 2023111014 \And
Yajat Lakhanpal \\ 2023111015 \And
Hrishiraj Mitra \\ 2023111037
}

\date{}

\begin{document}
\maketitle

\begin{abstract}
Graph Language Models (GLMs) condition answers on an input graph $G$, which
existing work treats as fixed. That breaks once the setting is conversational: if
a user revises a fact in language, the encoder still reads the old graph. We do
not propose a new architecture; we wrap an existing GLM in a loop that emits a
small edit, applies it to $G$, and re-encodes before the next question. Multi-turn
revision then produces a failure that loop cannot handle: a later utterance
incompatible with an edit already committed. We define such a contradiction as
unsatisfiability of the active constraint set, localize the earlier edit
responsible, and ask which revision to keep rather than preferring the newest. We
will evaluate on synthetic sessions where the gold resolution is deliberately
\emph{not} the most recent revision.
\end{abstract}

\section{Introduction}

A GLM answers questions from an input graph $G$, usually as a single pair
$(G, q)$. In a conversation a user may correct a fact: the language model sees
the correction, but the graph encoder still reads the original $G$, so later
answers can be grounded in a stale graph.

\citet{plenz-frank-2024-graph} introduced GLMs so structure and text are encoded
together. \citet{he-2024-gretriever} describe a conversational interface, but
GraphQA is still single-shot. TEA-GLM \citep{wang-2024-teaglm}, GraphToken
\citep{perozzi-2024-graphtoken}, and CLEGR \citep{petkar-2026-clegr} likewise
treat $G$ as fixed, and agentic graph token reasoning
\citep{peng-yang-2026-agentic} only re-chooses a view of it. GraphRAG keeps state
in an external store and SKG-Eval \citep{shil-samui-2026-skgeval} scores dialogue
against a parallel KG, but graph tokens sit inside the forward pass: neither
mutates the tensors the GNN encodes.

Edits are also cumulative: turn three applies to $G_2$, never to the original
graph. That introduces a second failure, because revision dialogue routinely
produces a later utterance that \emph{cannot} sit on top of an earlier one. If
turn~2 closes station~B and turn~5 puts through-service on the B--C line, turn~5
is neither a no-op nor a clean reversal, yet a write path with no notion of
conflict silently takes the newer edit and drops the closure. We propose a write
path that knows when a revision contradicts an earlier one, and asks.

\section{What Counts as a Contradiction}

``Contradiction'' is overloaded. We take it from belief revision and ontology
debugging, where a sentence set is inconsistent when it has no model, and ground
it on the graph. Let each revision $r$ denote constraints
$I(r)$ over attributed graphs --- the implied state of the edit, not the wording
of the utterance --- let $\mathrm{Inv}$ be the stated invariants (schema plus task
invariants such as connectivity or mutual exclusion), and let $R_{\mathrm{act}}$
be the revisions not yet superseded, so that
$T = \bigcup_{r \in R_{\mathrm{act}}} I(r) \cup \mathrm{Inv}$. A \emph{conflict}
holds iff $T$ has no model, and the \emph{culprit set} is a minimal
unsatisfiable subset of $T$. A conflict is a \textbf{contradiction} while no
culprit has been retired, and a \textbf{revision} if the user intends the new $r$
to supersede some $r_j$ in the culprit set.

This makes the DECODE judgment \citep{nie-2021-decode} decidable, and it is
inter-context conflict \citep{xu-2024-knowledge} with committed edits in place of
retrieved passages; localization corresponds to DECODE's
evidence indices, as in ontology debugging
\citep{shchekotykhin-2012-interactive}. Crucially, unsatisfiability is the
conflict, but whether it is a contradiction to surface or a revision to absorb
depends on which constraints remain \emph{active}, which is derivable neither
from $G$ nor from recency. That is what we propose to ask. Three cases arise:
\textbf{invariant violations}, caught by the schema check; \textbf{silent
incompatibilities}, solver-decidable given $\mathrm{Inv}$; and \textbf{intent
conflicts}, valid graphs that still cannot both be what the user meant.

\section{Method}

A GLM forward pass cannot rewrite $G$ inside the GNN, so we wrap an existing model
(G-Retriever, GraphToken, or TEA-GLM) in an outer loop. On each turn we prompt the
GLM with the current graph and $u_t$ to emit an edit $\Delta$, parse it (illegal
edits become no-ops), apply it with a schema-checked function, then re-encode $G$
and answer $q_t$.

The model has no tools. It already emits text, and we restrict that text to four
operations: \texttt{SET} a node attribute, \texttt{ADD} an edge, \texttt{DEL} an
edge, or \texttt{NOOP}. Two calls, same weights, and the QA prompt is only
$(G, q_t)$ --- the revision utterance is withheld, so if QA succeeds the fact is
genuinely in the graph. Constrained decoding or light LoRA is a fallback only if
writing the edit is the bottleneck. Skipping re-encoding makes this GraphRAG: it
is a GLM method only if the apply step precedes the next GNN forward pass.

Around this loop we add four stages. \textbf{Detect}: a conflict gate before
commit, using the schema check for invariant violations and a tentative apply plus
invariant test for silent ones. \textbf{Attribute}: on a conflict at turn $k$,
replay $G_0 \ldots G_{k-1}$ and isolate the minimal earlier edit set that makes
$u_k$ unsatisfiable --- a decision procedure over a typed log, not a learned
component. \textbf{Ask}: one discriminating question naming that edit.
\textbf{Repair}: retire the loser, replay the later edits that still apply, and
re-encode.

Detection and localization are machinery, not the contribution: every strategy we
know of resolves by recency, instruction-tuned models prefer the latest
instruction \citep{zhang-2025-iheval}, and models detect conflicting instructions
yet rarely disclose them \citep{he-2026-coninstruct}. Asking earns its place on
items where \textbf{recency is wrong} and the earlier edit should be preserved.

\section{Dataset}

Items come from a synthetic graph in the style of CLEGR: a fictional metro with
invented station names and attributes, so parametric memory cannot skip $G$. The
generator samples a graph $G_0$ of 15 to 30 nodes with
attributes such as open or closed, and two paths between a query pair so that
closing a node changes the answer. It fixes an executable query --- say the
shortest path through open stations --- samples \texttt{SET},
\texttt{ADD}, and \texttt{DEL} operations that change that answer, and verbalizes
each as an utterance. Gold answers come from a solver. We discard an item if the
question is answerable without the graph, if the utterance plus question recover
the answer, or if $G_0$ already yields the new answer.

For contradictions the generator must stop drawing every edit from the current
graph. To plant a conflict at turn $k$ against turn $j < k$, we sample the later
edit from an ancestor state $G_{j-1}$, so it is coherent alone but incompatible
with what the log has committed. Each item carries the
conflicting turn index, the gold culprit set, the intended graph, the stated
invariant where needed, and a \texttt{keep\_decision} in $\{$earlier, later,
amend$\}$. Gold answers stay solver-computable \emph{after} the decision: replay
the repaired log and solve on the intended graph. Only \texttt{keep\_decision} is
planted rather than derived, the one label neither the graph nor the log supplies;
simulated user replies come from it, so no human is needed at evaluation. The \textbf{recency-wrong subset} is the items where
\texttt{keep\_decision} is not \emph{later}; without it the rest is not worth
generating.

\section{Evaluation}

A frozen graph should keep the old answer after a revision; if it does not, the
item is too easy. Applying the gold edit and re-encoding, then asking $q_t$ with
no revision text, checks the GLM can read an updated graph; if prompting the
revision in text over a frozen $G_0$ matches that oracle, rewriting $G$ is
unnecessary. We compare the self-edit loop against a GraphRAG-style baseline that
never re-encodes.

We add four contradiction baselines, reported even where they should win. A
\textbf{symbolic solver} over the edit log replays, checks invariants, and
enumerates culprit sets, and should be near-exact on the first two cases; we
include it to make clear we claim neither detection nor localization as a result. A
\textbf{transcript-only LLM} reads a frozen graph plus the history;
\textbf{last-write-wins} commits the newest edit and fails by construction on the
recency-wrong subset; \textbf{oracle attribution} supplies the true culprit set,
so if accuracy does not move, localization is not the bottleneck. Beyond exact
match and edit correctness we report detection, localization, questions asked, and
post-repair state exactness against the intended graph, and that exactness on the
recency-wrong subset --- the decisive number. One risk: if contradictions are
imperceptible through graph tokens, detection falls to the symbolic gate; we will
probe this with an oracle-injected violation.

\bibliography{custom}

\end{document}
